%!TEX root = 00_main.tex
% ============================================================================
% CONTENT SLIDES
% ============================================================================
% Add your presentation content here.
% Create additional files (02_slides.tex, etc.) for more sections.
% ============================================================================

\section{Introduction}

% --- Motivation Slide ---
\begin{frame}{Wind Energy: Global Context}
    \begin{itemize}
        \item Wind power capacity exceeded 900 GW globally in 2023
        \item Offshore wind growing at 25\% annually
            \begin{itemize}
                \item Floating foundations enable deeper water installations
                \item Larger rotors (15+ MW turbines) improve capacity factors
            \end{itemize}
        \item Key engineering challenges:
            \begin{itemize}
                \item Fatigue life prediction under stochastic loading
                \item Aeroelastic stability of flexible blades
            \end{itemize}
        \item Cost reduction target: below 50 EUR/MWh by 2030
    \end{itemize}
\end{frame}

\section{Theory}

% --- Aerodynamics Slide ---
\begin{frame}{Wind Turbine Aerodynamics}
    The power extracted by a wind turbine is given by:
    \begin{equation}
        P = \frac{1}{2} \rho A C_p v^3
        \label{eq:power}
    \end{equation}
    where $\rho$ is air density, $A$ is rotor area, and $v$ is wind speed.

    \vspace{0.3cm}

    The Betz limit defines the maximum power coefficient:
    \begin{align}
        C_{p,\text{max}} &= \frac{16}{27} \approx 0.593 \label{eq:betz}\\
        a_{\text{opt}} &= \frac{1}{3} \label{eq:induction}
    \end{align}

    \vspace{0.3cm}

    Modern turbines achieve $C_p \approx 0.45$ (see Equation~\eqref{eq:betz}).
\end{frame}

\section{Results}

% --- Single Figure Slide ---
\begin{frame}{Power Curve Comparison}
    \begin{center}
        \includegraphics[height=4cm]{example_figure.png}
    \end{center}
    \begin{center}
        \textbf{Figure 1:} Measured vs. simulated power curve for a 5 MW offshore turbine.
    \end{center}
\end{frame}

% --- Two Figures Side-by-Side Slide ---
\begin{frame}{Blade Load Distribution}
    \begin{columns}[T]
        \begin{column}{0.4\textwidth}
            \centering
            \includegraphics[width=\textwidth]{example_figure.png}

            \textbf{(a)} Flapwise bending moment
        \end{column}
        \begin{column}{0.4\textwidth}
            \centering
            \includegraphics[width=\textwidth]{example_figure.png}

            \textbf{(b)} Edgewise bending moment
        \end{column}
    \end{columns}
\end{frame}

% --- Figure with Bullet Points Slide ---
\begin{frame}{Fatigue Damage Analysis}
    \begin{columns}[T]
        \begin{column}{0.45\textwidth}
            \begin{itemize}
                \item Rainflow counting applied to load time series
                \item Damage equivalent load (DEL): 12.5 MNm
                \item Critical location: blade root transition
                \item Safety factor of 1.35 applied per IEC 61400-1
                \item 20-year design life validated
            \end{itemize}
        \end{column}
        \begin{column}{0.45\textwidth}
            \centering
            \includegraphics[width=\textwidth]{example_figure.png}

            {\small\textbf{Figure:} S-N curve with Goodman correction}
        \end{column}
    \end{columns}
\end{frame}

\section{Conclusion}

% --- Summary Slide ---
\begin{frame}{Summary and Outlook}
    \begin{enumerate}
        \item Aeroelastic model validated against field measurements
        \item Fatigue life meets 20-year design requirements
        \item Blade pitch control reduces extreme loads by 15\%
        \item Future work: investigate wake effects in wind farm layout
    \end{enumerate}

    \vspace{0.5cm}

    \textbf{Thank you for your attention!}

    \vspace{0.3cm}

    Questions?
\end{frame}
